\par\noindent\textbf{(1)} 在芯包圆柱壁，轴向传播射线与光纤轴夹角为 \(\theta_1\)，其壁面入射角为 \(\pi/2-\theta_1\)。全反射要求 \(\sin(\pi/2-\theta_1)\ge n_2/n_1\)，即 \(\cos\theta_1\ge n_2/n_1\)，故 \(\sin\theta_1\le\sqrt{n_1^2-n_2^2}/n_1\)。
\par\noindent\textbf{(2)} 端面Snell定律 \(n_0\sin\theta_0=n_1\sin\theta_1\)。因此最大入射角满足 \(n_0\sin\theta_{0,\max}=\sqrt{n_1^2-n_2^2}\)，定义数值孔径 \(\mathrm{NA}=n_0\sin\theta_{0,\max}=\sqrt{n_1^2-n_2^2}\)，若右侧超过 \(n_0\) 时接收锥受外介质物理角限制。
\par\noindent\textbf{(3)} 定义相对折射率差 \(\Delta=(n_1-n_2)/n_1\ll1\)。则 \(n_1^2-n_2^2=n_1^2[1-(1-\Delta)^2]\approx2n_1^2\Delta\)，所以 \(\mathrm{NA}\approx n_1\sqrt{2\Delta}\)。该弱导展开对 \(\Delta^2\) 项的舍弃需明确。
\par\noindent\textbf{(4)} 真正光纤模结构取决于无量纲 \(V=(2\pi a/\lambda)\mathrm{NA}\)（\(a\) 为芯半径）。弱导阶跃光纤 \(V<2.405\) 时只有基本空间模式，而几何光学接收锥本身不保证多模或单模；靠近截止应求Maxwell本征方程。
